% ===================== RESUMEN EJECUTIVO =====================
\newpage
\section*{Resumen Ejecutivo}
\addcontentsline{toc}{section}{Resumen Ejecutivo}

El laboratorio objeto de estudio requiere un sistema de ventilación que garantice 12 renovaciones de aire por hora y una presión positiva controlada respecto a los espacios circundantes. La presurización positiva se logra mediante la inyección de aire limpio y filtrado a un caudal superior al de extracción, de modo que el flujo neto se dirija desde el interior hacia el exterior a través de aberturas controladas.

La solución de diseño propuesta utiliza un ventilador de impulsión directa ubicado en una pared del recinto, sin ductos de distribución. El aire ingresa al laboratorio a través de una boca circular equivalente a 412 mm de diámetro (radio 206.0 mm) y a una velocidad de 8 m/s. La salida del aire se efectúa mediante tres rejillas de exfiltración rectangulares de 353 mm $\times$ 336 mm cada una, operando a una velocidad facial de 3.0 m/s.

Los resultados principales indican un caudal de diseño de 64.0 m$^3$/min (3 840 m$^3$/h, 2 260 CFM), una potencia teórica del ventilador de 0.338 kW a las condiciones del sitio (Cajicá, $\rho$ = 0.88 kg/m$^3$) y una potencia instalada recomendada de 1.0 HP. El diseño cumple con los rangos de renovaciones de aire establecidos por ASHRAE 170, WHO y NIH para laboratorios de control medio y bioseguridad nivel 2. Los datos de contorno obtenidos son aplicables de manera directa a un modelo de dinámica de fluidos computacional (CFD) para la verificación del campo de presiones y la distribución de velocidades dentro del recinto.

\newpage
